\section{Implementation}
\label{sec:implementation}

mindtwin is implemented as an open-source Python package (MIT licence). All components run in a
single process without external services: twin state and history are kept in memory, the
knowledge graph uses an in-memory RDF store with SPARQL support, and episodic memory will use an
embedded SQLite database. This keeps every experiment reproducible from a plain package
installation, at the cost of horizontal scalability (Section~\ref{sec:discussion}).

\subsection{Knowledge layer}
\label{sec:impl-knowledge}

The knowledge layer consists of two OWL ontologies serialised in Turtle. The \emph{core} ontology
defines domain-independent concepts of a cognitive twin: assets and components (as
\texttt{sosa:FeatureOfInterest}), sensors and control signals (as \texttt{sosa:Sensor}), observable
properties, fault modes and events (anomalies, prognoses, maintenance actions). Structural
relations capture flow order between components (\texttt{upstreamOf}, transitive), mechanical
coupling (\texttt{mechanicallyCoupledTo}, symmetric) and sensor location (\texttt{monitors}).
Reusing SOSA~\citep{w3c2017sosa} keeps the graph readable by standard tools.

The \emph{turbofan} ontology specialises the core for the C-MAPSS engine: ten components (fan,
low- and high-pressure compressors and turbines, combustor, nozzle, bypass duct and two shafts),
the core and bypass gas paths, the two shafts, the location of all 21 recorded signals, and the
component degraded by each of the two fault modes. Only knowledge stated in the engine
documentation~\citep{saxena2008damage} is asserted by hand. Which sensors respond to which fault
(\texttt{hasSymptom}) is deliberately left empty: the twin learns these links from data and writes
them back as semantic memory, which keeps prior knowledge and learned knowledge separable for the
evaluation of RQ3.

\subsection{Twin core}
\label{sec:impl-core}

An \texttt{Asset} is built from a template individual in the ontology, so the twin and the
knowledge graph share one source of truth. It exposes structural queries used later by the
reasoning layer, such as the components downstream of a given component or the sensors that
monitor it. Notably, the five sensors whose correlation with RUL disappears when the fan fault is
present (Figure~\ref{fig:fault-signatures}) all monitor components downstream of the fan, which
illustrates how structure constrains fault signatures.

A \texttt{DigitalTwin} wraps one asset. It ingests one validated \texttt{Snapshot} per operating
cycle (operating settings and sensor readings), rejecting snapshots for another asset, out-of-order
cycles, unknown sensors and infinite values, and keeps the current state and the history (optionally
bounded). Cognitive components subscribe to the twin and are notified after every update; a
failing subscriber is logged and isolated so it cannot corrupt the twin's state. Ground-truth labels
such as RUL are never passed to the twin, so components cannot accidentally access them.

For experiments, each C-MAPSS engine is replayed cycle by cycle into its own twin. On a laptop
(Python~3.12, single core), the 249 training engines of FD004 (61{,}249 cycles) are replayed in
\SI{2.9}{\second}, about 21{,}000 snapshots per second, which is far faster than real time and
sufficient for all experiments. \todo{Report final hardware and re-measure at submission time.}

\todo{Cognition, memory and agent subsections.}
